% ============================================================
%  CAPÍTULO 7 — CONCLUSIONES
%  Memoria TFG: Transcripción automática de música de piano
%  Autor: María Piedad Castex Sierra
% ============================================================

\chapter{Conclusiones}
\label{cap:conclusiones}

El presente capítulo sintetiza el trabajo realizado, enumera las
contribuciones concretas del proyecto, reconoce sus limitaciones y
propone líneas de investigación y desarrollo para continuar el trabajo
iniciado.

% -----------------------------------------------------------
\section{Resumen del trabajo}
\label{sec:resumen_trabajo}
% -----------------------------------------------------------

\textbf{Problema y motivación.} La AMT de piano exige convertir audio
en una representación simbólica precisa: qué notas suenan, cuándo
empiezan y terminan, y con qué intensidad. Su naturaleza polifónica la
convierte en una tarea especialmente compleja, con aplicaciones de
interés en la educación musical, el análisis musicológico y la
preservación del patrimonio sonoro
\cite{benetos2019automatic, muller2021fundamentals}. Este trabajo
aborda el problema desde una perspectiva poco explorada en la
literatura: formular la transcripción como traducción de secuencia a
secuencia, en lugar de clasificación por tramas.

\textbf{Sistema desarrollado.} Se implementó \textit{Magdalena v1.0},
un sistema \textit{Encoder-Decoder} asimétrico entrenado sobre
MAESTRO:

\begin{itemize}
    \item \textbf{Codificador} (\textit{AudioSparsifinerEncoder}):
    adapta la atención dispersa de Wei et al. \newline\
    \cite{wei2023sparsifiner} a audio de longitud variable
    (Sección~\ref{sec:contribuciones}).
    \item \textbf{Decodificador} (\textit{MidiDecoder}):
    \texttt{nn.TransformerDecoder} estándar, vocabulario propio de
    311 \textit{tokens} \cite{oore2020time}.
    \item \textbf{Tres experimentos}: \texttt{v1\_baseline} (corpus
    completo), \texttt{exp01\_reduced\_data} (150 piezas) y
    \texttt{exp02\_short\_clips} (corpus completo, clips de la mitad
    de duración).
\end{itemize}

Ante la ausencia de un registro exportable de pérdida de validación,
el \textit{checkpoint} evaluado se seleccionó de forma manual mediante
el ranking interno de Colaboratory \newline(\texttt{v1\_baseline},
\texttt{exp02\_short\_clips}) o, cuando este no estuvo disponible,
mediante la última época (\texttt{exp01\_reduced\_data}).
\clearpage

\textbf{Resultados.} Sobre el conjunto de prueba
(Tabla~\ref{tab:resultados_test}):

\begin{itemize}
    \item F1 \textit{onset}: 0.103 (\texttt{exp02}) $>$ 0.080
    (\texttt{v1\_baseline}) $>$ 0.065 (\texttt{exp01}); ningún modelo
    supera F1 = 0.5 en ninguna pieza.
    \item Caída sistemática de F1 al añadir \textit{offset}
    (0.665--0.809): el sistema localiza bien el \textit{onset} pero
    apenas modela la duración de la nota.
    \item Parte de la dificultad es intrínseca a cada pieza y se
    comparte entre configuraciones; el desajuste entre notas
    predichas y de referencia correlaciona negativamente con el F1 en
    los tres modelos.
    \item \texttt{exp01\_reduced\_data} confirma que 150 muestras son
    insuficientes, agravado por ser el único experimento sin criterio
    de selección de \textit{checkpoint}.
    \item La convergencia temprana de \texttt{exp02\_short\_clips} no
    muestra evidencia de sobreajuste en las curvas registradas.
\end{itemize}

En conjunto, los resultados no compiten con el estado del arte --algo
esperable dadas las restricciones de escala y cómputo-- pero validan
la viabilidad de la arquitectura Sparsifiner en el dominio del audio
musical, y delimitan con precisión los factores sobre los que debe
incidir el trabajo futuro.
\clearpage
% -----------------------------------------------------------
\section{Contribuciones}
\label{sec:contribuciones}
% -----------------------------------------------------------

El presente proyecto ha realizado tres contribuciones diferenciadas al
campo de la transcripción automática de música para piano.

La primera es la \textbf{formulación integral de la AMT de piano como
un problema de traducción de secuencia a secuencia}. A diferencia de
los enfoques clásicos basados en clasificación por tramas que tratan
el tono, la duración y la dinámica de forma desacoplada, el sistema
desarrollado establece un mapeo directo y conjunto entre la
representación acústica continua y una secuencia simbólica de eventos
musicales. Para ello se diseñó un vocabulario propio de 311
\textit{tokens} que codifica de manera integral los cuatro atributos
fundamentales de cada nota (tono, inicio, duración y dinámica),
siguiendo la filosofía de representación por eventos de Oore et al.\
\cite{oore2020time} con adaptaciones específicas en la cuantización de
la velocidad (32 \textit{bins}) y en la resolución temporal
(pasos de 10\,ms).

La segunda contribución es la \textbf{adaptación del mecanismo de
atención dispersa Sparsifiner al dominio del audio de longitud
variable}. La arquitectura original de Wei et al.\
\cite{wei2023sparsifiner} fue diseñada para procesar imágenes de
dimensiones fijas, lo que la hace incompatible directamente con
secuencias de audio de duración arbitraria. Para superar esta
limitación se introdujeron tres modificaciones de ingeniería sobre la
implementación original:

\begin{itemize}
    \item \textbf{Recorte dinámico de la matriz de proyección temporal}
    (\textit{proj\_n}): en lugar de mantener una matriz de parámetros
    de tamaño fijo $\mathbf{P}_N \in \mathbb{R}^{T_{max} \times k}$
    y aplicarla sobre secuencias rellenas con \textit{padding}, se
    recorta dinámicamente a la longitud real de cada secuencia en cada
    paso de inferencia, eliminando el coste computacional y de memoria
    asociado al relleno.

    \item \textbf{Ajuste dinámico del presupuesto de atención}:
    la operación \newline \texttt{real\_k = min(self.attn\_budget, N)} garantiza
    que el presupuesto de conexiones activas nunca supere la longitud
    real de la secuencia, previniendo errores de índice en secuencias
    más cortas que la longitud máxima $T_{max}$.

    \item \textbf{Eliminación de la restricción de longitud fija}:
    se reemplazó la comprobación de longitud exacta por una cota superior la comprobación. La condición \texttt{assert num\_tokens == N} del código original, forzaba una longitud de entrada fija e
    incompatible con el procesamiento de fragmentos de audio de
    duración variable.
\end{itemize}

Estas tres modificaciones, tomadas en conjunto, permiten que el
codificador procese secuencias de cualquier longitud hasta $T_{max}$
sin necesidad de relleno hasta esa cota, lo que resulta especialmente
relevante para un corpus de duración variable como MAESTRO.

La tercera contribución es el \textbf{desarrollo de un sistema de
transcripción funcional de extremo a extremo}, incluyendo el
\textit{pipeline} completo de datos (preprocesamiento acústico,
tokenización MIDI y carga por lotes), el ciclo de entrenamiento con
precisión mixta y reanudación automática desde \textit{checkpoint}, y
el protocolo de evaluación con \texttt{mir\_eval} sobre el conjunto de
prueba oficial de MAESTRO. Este sistema constituye una base empírica
reproducible sobre la que contrastar el impacto del mecanismo
Sparsifiner frente a la atención densa estándar en el dominio del audio
musical.
\clearpage

% -----------------------------------------------------------
\section{Competencias desarrolladas}
\label{sec:competencias}
% -----------------------------------------------------------

En cuanto a las competencias enumeradas en la sección \ref{sec:objetivos} referentes a la especialidad de Computación, han sido trabajas en los siguientes aspectos del proyeto desarrollado:

\textbf{[CM3].} El diseño \textit{Encoder-Decoder asimétrico} de
\textit{Magdalena v1.0} es en sí mismo una decisión motivada por un
análisis de complejidad: la atención estándar tiene un coste
$O(N^2)$ en la longitud de la secuencia, inasumible para secuencias
de audio de hasta 70\,000 tramas, pero perfectamente manejable para
las secuencias de \textit{tokens} MIDI, mucho más cortas, que procesa
el decodificador. Esto llevó a reservar la atención dispersa
Sparsifiner --con coste sublineal en el número de conexiones
activas-- exclusivamente para el codificador, y a mantener atención
densa estándar donde la longitud de secuencia no lo justifica. Sobre
esa estrategia se implementaron además las tres adaptaciones de
ingeniería descritas en la Sección~\ref{sec:contribuciones} (recorte
dinámico de \textit{proj\_n}, ajuste dinámico del presupuesto de
atención y eliminación de la restricción de longitud fija), todas
ellas orientadas a garantizar el rendimiento del sistema dentro de las
restricciones de memoria de la GPU disponible (Sección~\ref{sec:limitaciones}).

\textbf{[CM4].} El proyecto recorre el ciclo completo de análisis,
diseño y construcción de un sistema inteligente: desde la revisión de
los fundamentos del Transformer y los paradigmas de traducción de
secuencia a secuencia (Capítulo~\ref{cap:estado_del_arte}), pasando
por el diseño razonado de la arquitectura \textit{Magdalena v1.0}
(Capítulo~\ref{cap:arquitectura}) --incluyendo decisiones como el uso
de codificación posicional aprendida en el codificador, motivada por
la naturaleza bidimensional del espectrograma siguiendo a Dosovitskiy
et al.\ \cite{dosovitskiy2021image}, frente a codificación sinusoidal
fija en el decodificador para favorecer la generalización a longitudes
no vistas--, hasta su implementación funcional de extremo a extremo y
su evaluación empírica (Capítulo~\ref{cap:resultados}).

\textbf{[CM5].} El vocabulario propio de 311 \textit{tokens}
(Sección~\ref{sec:contribuciones}) constituye una formalización
explícita de conocimiento musical humano --tono, \textit{onset},
\textit{offset} y dinámica de cada nota-- en una representación
discreta y computable, adaptando la filosofía de representación por
eventos de Oore et al.\ \cite{oore2020time} mediante decisiones de
cuantización propias (32 \textit{bins} de velocidad, resolución
temporal de 10\,ms). En el otro extremo del sistema, la
representación de espectrograma Log-Mel formaliza de forma análoga la
señal acústica en bruto en una representación computable adecuada
para el aprendizaje, cubriendo así el aspecto de \textit{percepción}
de la competencia. El aspecto de \textit{actuación} se aborda de
forma más limitada, al tratarse de una tarea de transcripción
--percepción hacia símbolo-- y no de un agente que opere sobre un
entorno; la generación autorregresiva de la secuencia de salida es, en
todo caso, la forma en que el sistema materializa ese conocimiento
formalizado en una decisión computable, símbolo a símbolo.

\textbf{[CM7].} Todo el proyecto gira en torno al diseño, desarrollo y
evaluación de un sistema de aprendizaje computacional: el
\textit{pipeline} de entrenamiento con precisión mixta, \textit{teacher
forcing} y reanudación automática desde \textit{checkpoint}
(Sección~\ref{sec:contribuciones}); el diseño de tres experimentos
controlados que varían sistemáticamente el volumen de datos y la
longitud de secuencia (\texttt{v1\_baseline},
\texttt{exp01\_reduced\_data}, \texttt{exp02\_short\_clips}); y un
protocolo de evaluación cuantitativa con \texttt{mir\_eval}
(Capítulo~\ref{cap:resultados}). El corpus MAESTRO, con
aproximadamente 200 horas de grabaciones, constituye precisamente el
gran volumen de datos del que el sistema extrae de forma automática
--sin reglas ni heurísticas hechas a mano-- el conocimiento simbólico
sobre los eventos musicales que contiene.

% -----------------------------------------------------------
\section{Limitaciones}
\label{sec:limitaciones}
% -----------------------------------------------------------

El sistema desarrollado presenta las siguientes limitaciones:

\begin{itemize}
    \item \textbf{Restricciones de hardware.} Entrenamiento en Google
    Colaboratory con una GPU T4 de 15\,GB, lo que impuso una
    configuración reducida ($d=256$, cuatro capas) frente a la de
    referencia del Transformer ($d=512$, seis capas), dificultando la
    comparación directa con el estado del arte.

    \item \textbf{Truncación de secuencia.} Audio limitado a 4\,096
    tramas ($\approx$82\,s) en \texttt{v1\_baseline} y
    \texttt{exp01\_reduced\_data}, y a 2\,048 tramas ($\approx$41\,s)
    en \texttt{exp02\_short\_clips}, descartando dependencias
    musicales a largo plazo en piezas más extensas.

    \item \textbf{Sin \textit{scheduler} de tasa de aprendizaje.}
    Tasa constante $\eta=10^{-4}$, sin \textit{warmup} ni decaimiento
    coseno \cite{vaswani2017attention}, pudiendo limitar la velocidad
    de convergencia y la calidad del mínimo alcanzado.

    \item \textbf{Sin aumento de datos.} Entrenamiento sobre el audio
    original de MAESTRO, sin \textit{pitch shifting},
    \textit{time stretching} ni ruido añadido, reduciendo la
    variabilidad efectiva del conjunto de entrenamiento.

    \item \textbf{\textit{Teacher forcing} sin \textit{curriculum
    learning}.} El uso estricto de \textit{teacher forcing} introduce
    \textit{exposure bias} entre entrenamiento e inferencia; no se
    implementaron mecanismos como \textit{scheduled sampling}.

    \item \textbf{DropPath inactivo.} Implementado en la arquitectura
    pero configurado con tasa cero (\texttt{drop\_path\_rate = 0.0})
    en todos los experimentos, actuando como función identidad.

    \item \textbf{Selección manual de \textit{checkpoint}.} Sin
    registro exportable de pérdida de validación, la selección se
    basó en el ranking interno no persistente de Colaboratory; en\newline
    \texttt{exp01\_reduced\_data}, no disponible, se empleó la última
    época sin ningún criterio de rendimiento, penalizando
    probablemente su resultado final. Las métricas F1 muestreadas
    periódicamente presentan además una variabilidad considerable en
    los tres experimentos.

    \item \textbf{Evaluación reducida y con ventanas no
    equivalentes.} Evaluación sobre 30 de las 125 piezas de prueba,
    sin metadatos por pieza (tempo, densidad polifónica, pedal).
    Además, la ventana de referencia de \texttt{exp02\_short\_clips}
    es aproximadamente la mitad de larga que la de los otros dos
    experimentos, por lo que los recuentos absolutos de notas --aunque
    no las métricas F1-- no son directamente comparables entre
    \texttt{exp02\_short\_clips} y el resto.

    \item \textbf{Especialización en piano.} Vocabulario de 88 notas
    del rango del piano estándar; la extensión a otros instrumentos
    requeriría rediseñar el vocabulario y reentrenar desde cero.
\end{itemize}

% -----------------------------------------------------------
\section{Trabajo futuro}
\label{sec:trabajo_futuro}
% -----------------------------------------------------------

Las limitaciones anteriores trazan las principales líneas de trabajo
futuro, ordenadas por impacto esperado sobre el rendimiento del
sistema:

\begin{itemize}
    \item \textbf{Configuración completa.} Entrenar con $d=512$,
    $H=8$, $L_{enc}=L_{dec}=6$ sobre hardware dedicado, permitiendo
    comparación directa con el estado del arte y midiendo cuánto de
    la brecha se debe a la restricción de escala.

    \item \textbf{Registro y selección automática de
    \textit{checkpoint}.} Exportar la pérdida de validación por época
    y automatizar la selección del mejor \textit{checkpoint} según
    una métrica objetivo (p.\,ej.\ F1 \textit{onset}), evitando el
    procedimiento manual actual.

    \item \textbf{\textit{Scheduler} de tasa de aprendizaje.}
    Incorporar calentamiento lineal y decaimiento coseno
    \cite{vaswani2017attention}, de alto impacto y bajo coste de
    implementación.

    \item \textbf{Aumento de datos.} Transposición tonal
    ($\pm 6$ semitonos), estiramiento temporal y ruido de fondo, para
    ampliar la variabilidad efectiva del corpus y mejorar la
    generalización.

    \item \textbf{Exploración sistemática de la longitud del clip.}
    Evaluar distintas longitudes de forma controlada y sobre ventanas
    de referencia equivalentes, dado que \newline\texttt{exp02\_short\_clips}
    obtuvo el mejor resultado global.

    \item \textbf{Decodificación avanzada.} \textit{Beam search},
    muestreo con núcleo o temperatura adaptativa, para reducir el
    \textit{exposure bias} de la decodificación voraz.

    \item \textbf{\textit{Chunking} solapado.} Sustituir la
    truncación dura por fragmentación con solapamiento, procesando
    piezas de duración arbitraria sin descartar información musical.

    \item \textbf{DropPath y ablación.} Explorar tasas no nulas
    ($p \in [0.1, 0.3]$) en configuraciones más profundas, junto con
    experimentos de ablación sobre cada componente del Sparsifiner.

    \item \textbf{Evaluación ampliada con metadatos.} Extender la
    evaluación a las 125 piezas del conjunto de prueba y registrar
    metadatos por pieza (tempo, densidad polifónica, pedal) para
    contrastar las hipótesis sobre la varianza de rendimiento
    observada.

    \item \textbf{Generalización a otros instrumentos.} Adaptar el
    vocabulario al rango tonal de instrumentos polifónicos como la
    guitarra o el bajo eléctrico, con un corpus equivalente a MAESTRO
    en calidad de alineamiento.

    \item \textbf{Interfaz gráfica y despliegue.} Desarrollo de la
    interfaz prevista en el diseño original, y optimización para
    inferencia en tiempo real mediante cuantización o destilación.
\end{itemize}